Modern machine learning services are increasingly deployed as distributed microservice systems, where resource allocation decisions must balance latency, reliability, and cost. In practice, however, autoscaling remains largely reactive: controllers add replicas only after the system is already under pressure, while safer deployments often rely on static over-allocation. Existing autoscalers do not learn a model of how scaling actions propagate through the service graph, and therefore cannot anticipate how a decision will affect future load, latency, cost, and service-level violations before it is executed. To fill this research gap, we present \emph{NetDream}, a graph world model for safe autoscaling of distributed machine learning services. NetDream learns action-conditioned dynamics over a microservice graph, predicting how load, latency, replicas, and violation risk evolve under candidate scaling actions. At runtime, NetDream uses this learned model for imagination-based planning: it rolls out multiple future scaling trajectories, filters out plans predicted to violate service-level constraints, and executes the first action of the best safe plan. This turns autoscaling from threshold-based reaction into topology-aware planning over a learned model of the system. We evaluate NetDream on a production-grade AWS Kubernetes deployment of Online Boutique, an 11-service, 14-edge microservice graph with a 66-dimensional system state and 48,000 collected transitions, under steady, variable, bursty, and flash-crowd traffic. NetDream-Safe is the most cost-efficient safe autoscaler among all evaluated methods. Each marginal replica-step it spends buys safety $2.7\times$ more cost-efficiently than the best static over-provisioning strategy. On variable-traffic workloads specifically, it achieves about $5\times$ fewer violations than standard reactive autoscaling while using only about half the cost of the strongest static over-provisioning baseline; on flash-crowd traffic the gain narrows because pod-startup delay bottlenecks all autoscalers. Finally, removing graph structure makes one-step prediction $11\times$ worse, showing that graph structure is required for reliable multi-step prediction. Code, trained models, infrastructure manifests, and a 30-second explainer video are released anonymously at \url{https://anonymous.4open.science/r/netdream-FCF9}.
